\section{Experimental Protocol}\label{sec:protocol}

\textbf{Metrics.} Classification tasks (1, 2, 3) use macro-F1 as the primary metric. Micro-F1
is also available for Task 1 and Task 2's main models, and AUC-PR for Task 1 and Task 3; Task
2's edge-free and 2-hop controls, evaluated for macro-F1 only, lack recorded
micro-F1/AUC-PR, so this paper does not report them for every model in every table -- only
where the evaluation produced them. Retrieval (Task 4) reports Recall@\{1,5,10\} in both
directions (audio$\to$text and text$\to$audio), plus a zero-shot tag-prediction check (tag
names embedded as text, no supervised training) as a sanity check that the learned space
carries semantic content at all.

\textbf{Model selection.} Every model evaluated on the held-out test set was selected by
best validation-set score (macro-F1 for classification, average Recall@10 for retrieval)
before that test evaluation. Task 3's edge-free \textit{gnn\_only} comparison
(Section~\ref{sec:results-t3}) is validation-only in a stricter sense: its training scripts
run a fixed 20 epochs and report the peak validation macro-F1 observed, rather than saving a
selected best checkpoint -- so "best" there describes the reported number, not a checkpoint
carried forward. This comparison is marked validation-only wherever reported, not presented
as a test-set result.

\textbf{Seeds.} Task 2/3's main comparisons in Sections~\ref{sec:results-t2}--\ref{sec:results-t3}
use a single training run each, matching the resources available; this is a real limitation,
stated as such rather than implied more robust than it is. The deployed GraphSAGE checkpoints
are measured against were trained with a documented seed (42); the Task 2 exact 2-hop control
script also fixes this same seed for its Python/NumPy/PyTorch state, so that comparison is
seed-matched. The edge-free set-baseline scripts (Task 2 and Task 3) do not set or record an
explicit seed, so those single-run contrasts are not seed-matched and are not guaranteed
exactly reproducible from a fresh run -- a further reason to read the edge-free numbers as
indicative rather than precise. Task 4's graph-vs-edge-free comparison, the paper's most
direct causal claim, is repeated across three training seeds and reported as mean $\pm$
standard deviation (Section~\ref{sec:results-t4}).

\textbf{Fairness controls.} Every graph-vs-edge-free comparison here holds node features,
optimizer, learning rate, batch size, epoch budget, and evaluation code fixed. One exception,
found on inspection rather than assumed away: Task 2's GraphSAGE training call does not pass
its configured dropout (0.4) through, so it silently trains at the model default of 0.2,
while the edge-free encoder is passed 0.4 explicitly -- not dropout-matched, on top of the
aggregation-operator difference below. We disclose this rather than silently correct it
(correcting would mean retraining GraphSAGE and changing every dependent Task 2 number);
Section~\ref{sec:cross-task} discusses its implications for the smallest Task 2 gap
specifically. Separately, replacing GraphSAGE layers with per-node MLPs removes the
aggregation operator itself, not topology alone, so this is an edge-free-encoder control, not
a pure edge-ablation. Neither guarantee extends to Task 2's non-learned 2-hop control, which
is not an edge-free comparison (it keeps GraphSAGE's topology) and uses a different training
recipe -- full-batch optimization, no dropout, a fixed 200-epoch budget, versus GraphSAGE's
mini-batches and early stopping (both share the same nominal weight decay, $10^{-4}$;
Section~\ref{sec:discussion} discloses this precisely and treats the 2-hop result as
exploratory, not a controlled ablation). Section~\ref{sec:results-t2} reports a cleaner
topology-only intervention (real vs.\ shuffled vs.\ random edges on the same architecture,
from a companion audit \cite{tismir_audit}) that isolates edge identity alone and finds the
same pattern. Parameter counts are reported for Task 2 (Table~\ref{tab:task2-full}); Task 3's
edge-free encoder is also smaller than its GraphSAGE counterpart (9,522 vs.\ 15,666
parameters, from the same companion evidence base \cite{tismir_audit}), though this paper's
own tables do not repeat that count for Task 3/4. Wherever reported, the edge-free
alternative uses fewer parameters than its graph counterpart (Section~\ref{sec:results}), so
no gap here can be explained by the edge-free model having more capacity.

\textbf{Class weighting.} Weighting was not held fixed across every Task 3 number, stated
precisely rather than implied otherwise. Table~\ref{tab:task3-full}'s \textit{bert\_only}
(0.7298), \textit{concat} (0.7061), and \textit{cross\_attention} (0.6951) rows are all
\emph{unweighted}, directly comparable to each other. The two \textit{gnn\_only} rows are a
before/after methodology pair, not comparable conditions: 0.0112 is the original unweighted
run that first exposed a label-imbalance problem, and 0.1621 is the same setup with
positive-class weighting added to fix it. The edge-free \textit{gnn\_only} control (0.1599
mean-pool, 0.1689 max-pool, Section~\ref{sec:results-t3}) uses the same positive-class
weighting, matching the 0.1621/0.1640-family comparators, not the unweighted 0.0112 row. A
positive-weighted \textit{cross\_attention} variant also exists separately (0.6769, lower
than the unweighted 0.6951 here) but is not used in this paper's comparisons -- not
selectively omitted, it simply answers a different question (does weighting change
\textit{cross\_attention} specifically). No positive-weighted \textit{bert\_only} or
\textit{concat} run exists. Task 1 does not use positive-class weighting.
